% ============================================================
%  6.3 Técnicas y algoritmos relacionados   (capítulo 6, Marco teórico)
%  Se incluye desde secciones/07_marco_teorico.tex con \input.
% ============================================================
\section{Técnicas y algoritmos relacionados}
\label{sec:tecnicas-algoritmos}

Las técnicas que integran la comparación se organizan como una escalera de complejidad
creciente, que va desde líneas base de tasas, sin parámetros aprendidos más allá de
promedios, hasta modelos de aprendizaje automático no lineales. Esta estrategia se
fundamenta en el teorema de \enquote{no hay almuerzo gratis} (\textit{no free lunch}) de
\textcite{wolpert1996}, según el cual, promediado sobre todos los problemas posibles, ningún
algoritmo de aprendizaje supera a otro, de modo que la superioridad de una técnica solo
puede establecerse empíricamente para una clase de problemas concreta. En consecuencia, la
práctica recomendada consiste en comparar los modelos complejos con líneas base simples e
interpretables: si un modelo sofisticado no mejora de manera sustantiva a una tasa
proporcional a la población, su complejidad adicional no está justificada. Cada peldaño de
la escalera responde a una hipótesis distinta sobre qué información explica las consultas
esperadas, y su inclusión se justifica en esos términos.

\subsection{Líneas base de tasas: tasa global y tasas estratificadas}
\label{subsec:lineas-base}

La línea base más simple supone que las consultas son proporcionales a la población, con
una tasa común a todo el territorio. La tasa global se estima como
$\hat{\lambda} = \sum_i Y_i / \sum_i P_i$, y las consultas esperadas de cada zona como
$\hat{\mu}_i = \hat{\lambda} \cdot P_i$. Este modelo encarna la hipótesis nula sustantiva
del proyecto, según la cual la población basta para explicar las consultas, y constituye por
ello el modelo de referencia contra el que se miden todos los demás (véase
\ref{subsec:modelo-referencia}).

Un primer refinamiento consiste en estratificar la tasa. Si $g(i)$ denota el estrato al que
pertenece la zona~$i$, la tasa del estrato se estima como
$\hat{\lambda}_g = \sum_{i \in g} Y_i / \sum_{i \in g} P_i$, y las consultas esperadas como
$\hat{\mu}_i = \hat{\lambda}_{g(i)} \cdot P_i$. Este procedimiento es análogo a la
estandarización de tasas por estratos empleada en demografía y epidemiología, y permite
incorporar heterogeneidad sin estimar modelos. En el proyecto se consideran dos
estratificaciones: por municipio, que captura diferencias institucionales, de adopción o de
cobertura cartográfica entre los municipios del eje metropolitano, y por anillos de
distancia al centro, que captura el gradiente centro-periferia.

La estratificación por anillos se fundamenta en la economía urbana, cuyos modelos clásicos
de localización, como el de \textcite{alonso1964}, predicen que la intensidad de uso del
suelo y la densidad disminuyen con la distancia al centro de negocios, y cuyos estudios
empíricos, desde \textcite{clark1951}, han descrito la densidad de población urbana como una
función aproximadamente exponencial negativa de esa distancia. Por analogía, cabe esperar
que la tasa de consultas por habitante también varíe con la distancia al centro, ya sea por
la concentración de actividades y destinos, por diferencias socioeconómicas en la adopción
de la aplicación o por la desigual densidad de rutas mapeadas. Las líneas base
estratificadas permiten examinar si esta estructura simple explica una parte sustancial de
la variación antes de recurrir a modelos más complejos.

\subsection{Suavizado espacial por vecindad}
\label{subsec:suavizado-vecindad}

El segundo peldaño aprovecha la autocorrelación espacial mediante un estimador de suavizado
local. Para cada zona~$i$, con primer anillo de vecinas $N(i)$, que excluye a la propia
zona, las consultas esperadas se estiman como
\[
  \hat{\mu}_i = P_i \cdot \frac{\sum_{j \in N(i)} Y_j}{\sum_{j \in N(i)} P_j},
\]
es decir, la población de la zona multiplicada por la tasa conjunta de sus vecinas. El
estimador encarna la hipótesis de que zonas próximas comparten las condiciones que
determinan la tasa de consultas, como la adopción, la cobertura cartográfica o el perfil
socioeconómico, en coherencia con la primera ley de Tobler.

Este estimador se relaciona con dos familias de métodos. Por una parte, con los estimadores
empírico-bayesianos de tasas: \textcite{clayton1987} y \textcite{marshall1991} propusieron
contraer la tasa bruta de cada área hacia una tasa de referencia, global o local, con un
peso que aumenta cuando la población del área es pequeña, y el suavizado local por vecindad
puede interpretarse como el caso límite en el que la tasa propia recibe peso nulo y la tasa
vecinal peso total; programas de análisis espacial como GeoDa han difundido estos suavizados
de tasas en la práctica aplicada \parencite{anselin2006}. Por otra parte, se relaciona con
la regresión por vecinos más cercanos, en la que la predicción de una unidad se obtiene como
promedio de las respuestas de las unidades más próximas; aquí la proximidad se define por
contigüidad hexagonal y el promedio se pondera por la población.

La aplicación de este estimador requiere cuidados específicos para evitar la fuga descrita
en \ref{subsec:fuga-datos}. La tasa vecinal debe calcularse sin incluir la propia zona y, en
cada iteración de validación, solo con las zonas del conjunto de entrenamiento; las zonas
cuyo primer anillo queda íntegramente fuera del entrenamiento, situación frecuente en el
interior de los bloques de validación, requieren una regla de respaldo declarada de
antemano, como recurrir al segundo anillo o a la tasa global. Estas precauciones garantizan
que el suavizado se evalúe en las mismas condiciones en las que se aplicaría a una zona cuyas
consultas no se conocen.

\subsection{Modelos lineales generalizados: regresión de Poisson con offset}
\label{subsec:glm-poisson}

Los modelos lineales generalizados (GLM), introducidos por \textcite{nelder1972} y
desarrollados sistemáticamente por \textcite{mccullagh1989}, extienden la regresión lineal a
variables respuesta de la familia exponencial mediante tres componentes: una distribución
para la respuesta, un predictor lineal $\eta_i = x_i^{\top}\beta$ y una función de enlace $g$
que relaciona la media con el predictor, $g(\mu_i) = \eta_i$. La regresión de Poisson
corresponde a la distribución de Poisson con enlace logarítmico. Incorporando la población
como exposición, el modelo se escribe $\log \mu_i = \log P_i + \beta_0 + x_i^{\top}\beta$, o
de manera equivalente $\mu_i = P_i \cdot \exp(\beta_0 + x_i^{\top}\beta)$, lo que en notación
compacta se expresa como $\log \mu = \log P + X\beta$. Los parámetros se estiman por máxima
verosimilitud mediante mínimos cuadrados reponderados iterativamente
\parencite{nelder1972}.

La principal ventaja de este peldaño es su interpretabilidad: cada coeficiente representa el
cambio en el logaritmo de la tasa de consultas por habitante asociado a una unidad de la
variable correspondiente, de modo que $\exp(\beta_k)$ es una razón de tasas. Por ejemplo, un
coeficiente asociado a la cobertura se interpretaría como el factor multiplicativo de la tasa
en las zonas cubiertas respecto de las no cubiertas, manteniendo constantes las demás
variables. En términos de implementación, la biblioteca statsmodels admite directamente un
\textit{offset} en sus modelos GLM, mientras que en scikit-learn, cuyo
\var{PoissonRegressor} no admite \textit{offset}, puede obtenerse un ajuste equivalente
modelando la tasa $y_i / P_i$ con pesos muestrales iguales a $P_i$. Este modelo encarna la
hipótesis de que las variables de contexto actúan de forma multiplicativa y aditiva en la
escala logarítmica, sin interacciones no especificadas.

\subsection{Regresión binomial negativa y modelos con exceso de ceros}
\label{subsec:binomial-negativa}

Cuando existe sobredispersión, la regresión binomial negativa ofrece una generalización
natural de la de Poisson. En su parametrización más común, denominada NB2, la media conserva
la misma estructura, $\log \mu_i = \log P_i + x_i^{\top}\beta$, pero la varianza pasa a ser
$\operatorname{Var}(Y_i) = \mu_i + \alpha \mu_i^2$, donde $\alpha \ge 0$ es el parámetro de
dispersión; cuando $\alpha$ tiende a cero, el modelo se reduce al de Poisson
\parencite{hilbe2011,cameron2013}. La binomial negativa puede derivarse como una mezcla de
Poisson con heterogeneidad no observada de distribución gamma, lo que le confiere una
interpretación sustantiva: zonas con las mismas características observadas pueden diferir en
factores no medidos, como la intensidad de difusión local de la aplicación. Su inclusión en
la comparación se justifica porque, ante sobredispersión, ofrece errores estándar más
realistas y pondera de forma distinta las zonas de conteos altos.

Cuando el exceso de ceros persiste incluso tras considerar la sobredispersión, pueden
emplearse modelos inflados de ceros, que combinan un proceso binario que genera ceros
estructurales con probabilidad~$\pi$ y un proceso de conteo, de modo que
$P(Y = 0) = \pi + (1 - \pi) \cdot f(0; \mu)$ \parencite{lambert1992}, o modelos de valla
(\textit{hurdle}), que separan la decisión de si hay o no eventos de la cantidad de eventos
cuando los hay. La decisión entre estas alternativas no debe basarse en un único contraste
mecánico, sino en la combinación de varios criterios: la comparación entre el número de
ceros observados y el esperado bajo el modelo binomial negativo, los criterios de
información y el desempeño en validación, y la plausibilidad sustantiva de un proceso
generador de ceros estructurales. Conviene recordar que la binomial negativa, por su mayor
varianza, predice más ceros que la Poisson con la misma media, por lo que con frecuencia
absorbe buena parte del aparente exceso de ceros.

\subsection{Gradient boosting con pérdida de Poisson}
\label{subsec:gradient-boosting}

El \textit{gradient boosting}, formalizado por \textcite{friedman2001}, construye un modelo
aditivo $F(x) = \sum_m \nu \cdot h_m(x)$ mediante la incorporación secuencial de aprendices
débiles, normalmente árboles de decisión poco profundos, cada uno ajustado al gradiente
negativo de la función de pérdida respecto de las predicciones del modelo acumulado, con una
tasa de aprendizaje~$\nu$ que controla la contribución de cada árbol. Su atractivo para el
proyecto reside en que captura relaciones no lineales e interacciones entre variables sin
necesidad de especificarlas de antemano, algo que un GLM solo logra si el analista las
introduce explícitamente. Cuando la pérdida es la devianza de Poisson y el enlace es
logarítmico, el modelo estima directamente el logaritmo del conteo esperado, conservando la
coherencia con la naturaleza de los datos.

Existen varias implementaciones eficientes. XGBoost \parencite{chen2016} introdujo un
objetivo regularizado, un algoritmo de búsqueda de cortes consciente de la dispersión de los
datos y optimizaciones de sistema que lo convirtieron en referencia en la práctica; ofrece
el objetivo \var{count:poisson} y permite incorporar la exposición como margen base
(\var{base_margin}) igual a $\log P_i$. LightGBM \parencite{ke2017} acelera el entrenamiento
mediante el muestreo de observaciones basado en gradientes y la agrupación de variables
mutuamente excluyentes, ofrece un objetivo de Poisson y admite la exposición como puntuación
inicial (\var{init_score}). En scikit-learn, \var{HistGradientBoostingRegressor}, basado en
histogramas, admite \var{loss="poisson"} desde la versión 0.23; según la documentación
oficial, esta pérdida implementa internamente la mitad de la devianza de Poisson con enlace
logarítmico y requiere $y \ge 0$ \parencite{sklearn-hgb}. Como este estimador no admite
\textit{offset}, la exposición puede incorporarse modelando la tasa $y_i / P_i$ con pesos
muestrales $P_i$, formulación equivalente en esperanza a la del \textit{offset}.

La flexibilidad del \textit{gradient boosting} implica un mayor riesgo de sobreajuste, por
lo que su uso requiere controlar hiperparámetros como el número de iteraciones, la
profundidad o el número de hojas de los árboles, la tasa de aprendizaje y la
regularización, idealmente mediante parada temprana y búsqueda de hiperparámetros dentro de
la propia validación cruzada espacial, sin tocar el conjunto de prueba. Su menor
interpretabilidad puede compensarse parcialmente con gráficos de dependencia parcial. Este
peldaño encarna la hipótesis de que la relación entre contexto y tasa de consultas es no
lineal o presenta interacciones relevantes, y su inclusión permite verificar empíricamente
si esa complejidad adicional se traduce en una mejora sobre modelos más simples.

\subsection{Preparación de variables geoespaciales}
\label{subsec:preparacion-geoespacial}

La construcción del conjunto de datos analítico requiere varias operaciones geoespaciales.
La primera es la asignación de cada consulta a una zona, que se realiza convirtiendo las
coordenadas de origen (latitud y longitud en WGS84) al identificador de la celda H3 de
resolución~8 que las contiene, mediante la función \var{latlng_to_cell} de la biblioteca H3
en su versión~4.x; la misma operación, aplicada a otras resoluciones, genera las variantes de
sensibilidad, y la función que devuelve la celda padre de resolución~6 asigna cada zona a su
macrozona. Previamente, las consultas se depuran de registros con coordenadas inválidas o
situadas fuera del área de estudio y de posibles repeticiones de la misma consulta por un
mismo identificador de instalación en intervalos muy breves, y se agregan los conteos por
zona.

La segunda operación es la integración con la población y con la oferta. Dado que Kontur
Population se publica sobre la misma grilla H3 de resolución~8, la unión con las zonas se
realiza por identificador de celda, sin pérdida por reproyección; para otras resoluciones,
la población se agrega sumando las celdas hijas o se reparte según corresponda. La cobertura
se calcula como la distancia mínima entre la geometría de la zona (o su centroide) y los
trazados del feed GTFS reconstruidos como líneas a partir de \var{shapes.txt}. Estas
distancias deben calcularse en un sistema de coordenadas proyectado métrico, como la
proyección UTM zona 19~Sur (EPSG:32719) que corresponde a Cochabamba, y no en grados
geográficos. Las uniones espaciales (\textit{spatial joins}) de GeoPandas permiten además
asignar a cada zona el municipio en el que se sitúa la mayor parte de su superficie.

La tercera operación es la construcción de las variables de contexto: la distancia del
centroide de la zona al centro de la ciudad y su discretización en anillos; el municipio; el
indicador de cobertura según el umbral de 500~m y la distancia continua al trazado; y las
variables de vecindad, como la población del primer anillo o la tasa vecinal calculada con
las precauciones contra la fuga descritas en \ref{subsec:fuga-datos}. Todas estas
transformaciones se implementan como funciones reproducibles que se ajustan, cuando tienen
parámetros dependientes de los datos, solo con las zonas de entrenamiento de cada
iteración.

\subsection{Entorno tecnológico}
\label{subsec:entorno-tecnologico}

El proyecto se desarrolla en Python, lenguaje que concentra el ecosistema de bibliotecas
abiertas más amplio para la Ciencia de Datos y el análisis geoespacial. La manipulación de
datos tabulares se realiza con pandas \parencite{mckinney2010}, y la de datos geoespaciales
vectoriales con GeoPandas \parencite{jordahl2020}, que extiende pandas con geometrías y
operaciones espaciales como las uniones y los cálculos de distancia. La indexación en la
grilla hexagonal se efectúa con h3-py, el enlace oficial para Python de la biblioteca H3
\parencite{h3-overview}. El análisis de autocorrelación espacial, incluidos el $I$ de Moran,
los indicadores LISA y la construcción de matrices de pesos, se realiza con PySAL y su
módulo esda \parencite{rey2007}.

Para el modelado se emplean scikit-learn \parencite{pedregosa2011}, que proporciona los
estimadores de \textit{gradient boosting} basados en histogramas, las utilidades de
validación cruzada por grupos y las métricas de devianza; statsmodels
\parencite{seabold2010}, que ofrece los modelos lineales generalizados de Poisson y binomial
negativa con \textit{offset} e inferencia estadística completa; y, de forma complementaria,
XGBoost \parencite{chen2016} y LightGBM \parencite{ke2017}. Los mapas interactivos de
estimaciones y brechas se generan con Folium, biblioteca que produce mapas web basados en
Leaflet \parencite{folium}. La reproducibilidad se asegura fijando las versiones de las
bibliotecas y las semillas aleatorias, y organizando el código como una secuencia de guiones
que puede reejecutarse de principio a fin.
